\scanpage{106}

\mspage{112}

entre les sous-catégories pleines formées des objets loc. constants.

$b'_2)$ Comme $b'_1)$, mais avec $\mathbb D$-équiv. $2$-\uline{complète} au lieu de $1$-complète, i.e. $f^*$ induit une équivalence entre les sous-catégories formées des coeff. loc. const.

$b''_{\mathrm{coh}})$ Si $f:Z'\to Z$ de $\mathrm{Diag}$ est une Hot-équivalence, alors c'est une $\mathbb D$-équivalence cohomologique complète, i.e. pour tout $F\in\mathrm{Ob}\,\mathbb D(Z)$ qui est \uline{loc. constant},

\[
\text{on a}\qquad H^{\unsure{\cdot|i}}_{\mathbb D}(Z,F)\xrightarrow{\sim}H^{\cdot}_{\mathbb D}(Z',f^{\unsure{2|*}}F).
\]

$b''_{\mathrm{hom}})$ Si $f:Z'\to Z$ de $\mathrm{Diag}$ est une Hot-équivalence, c'est une $\mathbb D$-équivalence homologique complète, i.e. pour tout $F\in\mathrm{Ob}\,\mathbb D(Z)$ qui est \uline{loc. constant}, on a

\[
H_{\cdot}^{\mathbb D}(Z',f^*(F))\xrightarrow{\sim}H_{\cdot}^{\mathbb D}(Z,\unsure{F|E}).
\]

% slanted
\margin{Cf variantes $\pm$ spéciales pour HOT, pages \unsure{116 à 125|116/125}}

On a le diagramme d'implications tautologiques pour les variantes de Der 3 :

\diagram{\begin{tikzcd}[ampersand replacement=\&,column sep=small,row sep=small]
 \& b'_2\arrow[d,Rightarrow]\arrow[drr,Rightarrow] \& \& {}\\
 \& b'_1\arrow[dl,Rightarrow]\arrow[dr,Rightarrow]\arrow[dd,Rightarrow] \& \& a'\arrow[ddd,Rightarrow]\\
 b''_{\mathrm{coh}}\arrow[dr,Rightarrow] \& \& b''_{\mathrm{hom}}\arrow[dl,Rightarrow] \& {}\\
 \& b\arrow[drr,Rightarrow] \& \& {}\\
 \& \& \& a
\end{tikzcd}}

Dans tous les cas à ma connaissance, les dérivateurs satisfont la condition la plus forte Der 3 $b'_2$.

% archive sheet number: 106
